% !TeX program = pdflatex
%% Copyright 2026 Petr Socha
%%
%% This work may be distributed and/or modified under the
%% conditions of the LaTeX Project Public License, either version 1.3c
%% of this license or (at your option) any later version.
%% The latest version of this license is in
%%   https://www.latex-project.org/lppl.txt
%% and version 1.3c is part of all distributions of LaTeX
%% version 2008 or later.
%%
%% This work has the LPPL maintenance status `maintained'.
%% The Current Maintainer of this work is Petr Socha.
%% Maintenance and bug reports: https://github.com/petrsocha/pslib-latex/issues
%%
%% The components of this work are listed in MANIFEST.txt.
%% The school logo is a separately licensed third-party asset.

% ============================================================================
% ODBORNÁ PRÁCE -- SPŠ A VOŠ LIBEREC
% ============================================================================
% Tento soubor slouží k nastavení práce a k výběru jejích částí.
% Vlastní text pište do souborů ve složce kapitoly/.
%
% JAK ZAČÍT:
%   1. Vyplňte údaje o práci a oba abstrakty níže.
%   2. Upravte čestné prohlášení a případně vložte PDF zadání.
%   3. Otevřete soubory ve složce kapitoly/ a pište svou práci.
%   4. V oddílu KAPITOLY můžete soubory přidávat, odebírat a přesouvat.
%
% Řádky začínající znakem % jsou komentáře a neobjeví se ve výsledném PDF.
% V Overleafu zvolte překladač pdfLaTeX. Soubor spslib.cls neupravujte.
%
% POZOR: Značky %TC:ignore a %TC:endignore slouží k počítání slov.
% Neměňte je ani nepřesouvejte; počítají se jen připojené kapitoly.
% ============================================================================

%TC:ignore
\documentclass{spslib}
\addbibresource{literatura.bib} % Zdroje pro automatické citace.

% ----------------------------------------------------------------------------
% 1. ÚDAJE O PRÁCI -- VYPLŇTE
% ----------------------------------------------------------------------------
% Mění se pouze text uvnitř složených závorek { }.
% Typ práce může být např. "Maturitní práce" nebo "Ročníková práce".
\typprace{Maturitní práce}
\nazevprace{Ukázková odborná práce}
% Volitelné: krátký název se používá pouze v záhlaví stránek.
% Je vhodný pro delší názvy práce. Pokud jej nezadáte, použije se plný název.
%\kratkynazevprace{Ukázková práce}
\autorprace{Jan Novák}
\oborprace{Informační technologie}
\vedouciprace{Ing. Jana Nováková}
\skolnirok{2026/2027}
\datumprohlaseni{17. října 2026} % Datum uvedené u podpisu prohlášení.

% ----------------------------------------------------------------------------
% 2. ANOTACE A ANGLICKÉ SHRNUTÍ -- VYPLŇTE
% ----------------------------------------------------------------------------
% Stručně představte téma, cíl, postup a výsledek své práce.
% Český text se vypíše pod nadpisem Anotace, anglický pod Summary.
% Pokud jsou delší, sazba automaticky pokračuje na další stránce.
\anotace{%
  Zde bude stručná anotace práce v češtině.
}
\summary{%
  A short summary of the project in English.
}

% ----------------------------------------------------------------------------
% 3. ČESTNÉ PROHLÁŠENÍ -- ZKONTROLUJTE A PODLE POTŘEBY UPRAVTE
% ----------------------------------------------------------------------------
% Upravujte pouze text prohlášení, nikoliv jeho nadpis a podpisové pole.
% Datum a jméno u podpisu se přebírají z údajů o práci výše.
\prohlaseni{%
  Prohlašuji, že jsem tuto práci vypracoval/a samostatně.
  V práci jsem použil/a pouze zdroje uvedené v seznamu literatury
  a všechny převzaté části, včetně textů, obrázků, kódu a výstupů
  umělé inteligence, jsou řádně označeny a citovány.
}

% ----------------------------------------------------------------------------
% 4. PODEPSANÉ ZADÁNÍ -- VOLITELNÉ
% ----------------------------------------------------------------------------
% Pokud máte naskenované zadání, vložte do projektu dvoustránkové PDF.
% První strana PDF je líc zadání, druhá rub (i kdyby byl prázdný).
% Odstraněním znaku % na následujícím řádku nahradíte zástupný list.
% Bez tohoto řádku zůstane v dokumentu místo pro fyzické vložení zadání.
%\zadani{zadani.pdf}

% ----------------------------------------------------------------------------
% 5. SESTAVENÍ ÚVODNÍCH STRAN -- PONECHTE BEZE ZMĚNY
% ----------------------------------------------------------------------------
\begin{document}
\prednistrany % Titulní list, zadání, anotace, prohlášení a obsah.

% ----------------------------------------------------------------------------
% 6. KAPITOLY -- PŘIDÁVEJTE, ODEBÍREJTE A PŘESKUPUJTE
% ----------------------------------------------------------------------------
% Každý příkaz \input vloží obsah jednoho souboru ze složky kapitoly/.
% Na obrázky, tabulky, kapitoly a výpisy odkazujte přes \autoref{...}.
% Automaticky doplní typ objektu (např. „obrázek 2.1“).
% Nepotřebnou kapitolu lze vypnout zakomentováním celého řádku znakem %.
% Novou kapitolu přidejte jako soubor .tex a vložte další příkaz \input.
% Číslování kapitol a odkazy se upravují automaticky.
% Úvod a Závěr jsou nečíslované, ale objeví se v obsahu.
%
% DŮLEŽITÉ: Následující seznam je jediné místo, podle kterého se počítají
% slova na titulní straně. Zachovejte kolem něj značky %TC.
%TC:endignore
%% Copyright 2026 Petr Socha
%%
%% This work may be distributed and/or modified under the
%% conditions of the LaTeX Project Public License, either version 1.3c
%% of this license or (at your option) any later version.
%% The latest version of this license is in
%%   https://www.latex-project.org/lppl.txt
%% and version 1.3c is part of all distributions of LaTeX
%% version 2008 or later.
%%
%% This work has the LPPL maintenance status `maintained'.
%% The Current Maintainer of this work is Petr Socha.
%% Maintenance and bug reports: https://github.com/petrsocha/pslib-latex/issues
%%
%% The components of this work are listed in MANIFEST.txt.
%% The school logo is a separately licensed third-party asset.

% ============================================================================
% ÚVOD -- ZÁKLADNÍ TEXT A ODSTAVCE
% ============================================================================
% Uvod je nečíslovaný, ale bude uvedeny v obsahu dokumentu.
% Běžně píšete text do odstavců. Nový odstavec vytvoříte prázdným řádkem.
\necislo{Úvod}

Tento dokument obsahuje jednoduché ukázky sazby odborné práce v~LaTeXu.
Jeho cílem je předvést nejběžnější postupy, které se při psaní práce používají.

První odstavec může stručně představit téma. V~dalším odstavci obvykle
následuje \emph{motivace}, \textbf{hlavní cíl} a~krátký popis postupu.
Jednotlivé ukázkové kapitoly lze později nahradit vlastním textem.

% Jednopísmenné předložky zapisujte s nezlomitelnou mezerou: v~textu, k~mereni.
% LaTeX ji vytvoříte znakem ~. Nadpisy a zalomení stran neměňte ručně.

%% Copyright 2026 Petr Socha
%%
%% This work may be distributed and/or modified under the
%% conditions of the LaTeX Project Public License, either version 1.3c
%% of this license or (at your option) any later version.
%% The latest version of this license is in
%%   https://www.latex-project.org/lppl.txt
%% and version 1.3c is part of all distributions of LaTeX
%% version 2008 or later.
%%
%% This work has the LPPL maintenance status `maintained'.
%% The Current Maintainer of this work is Petr Socha.
%% Maintenance and bug reports: https://github.com/petrsocha/pslib-latex/issues
%%
%% The components of this work are listed in MANIFEST.txt.
%% The school logo is a separately licensed third-party asset.

% ============================================================================
% ANALÝZA -- NADPISY, SEZNAMY A CITACE
% ============================================================================
\chapter{Analýza}
\label{chap:analyza}

% Nadpisy v LaTeXu: \section, \subsection a případně \subsubsection.
% Čísla nadpisů se doplňují automaticky. Nezapisujte je ručně.
\section{Práce s odbornými zdroji}

Základní principy analogových a~digitálních obvodů lze studovat například
v~knize \emph{Foundations of Analog and Digital Electronic Circuits}~\cite{agarwal2005}.

% Citaci vložte pomocí \cite{klic}. Klíč musí existovat v literatura.bib.
% Ukázka: odborná kniha (typ @book). Další druhy zdrojů najdete v ostatních kapitolách.
% Pokud čerpáte konkrétní myšlenku, umístěte citaci k příslušné větě.

\section{Požadavky}
\label{sec:pozadavky}

% Odrážkový seznam začíná prostředím itemize.
\subsection{Funkční požadavky}
Příklady požadavků na jednoduchou aplikaci:
\begin{itemize}
  \item zobrazit naměřenou hodnotu,
  \item uložit výsledek měření,
  \item upozornit na překročení limitu.
\end{itemize}

% Číslovaný seznam se zapisuje do prostředí enumerate.
\subsection{Nefunkční požadavky}
Pro ilustraci lze vymezit i další požadavky:
\begin{enumerate}
  \item snadná obsluha,
  \item přehledné rozhraní,
  \item spolehlivé ukládání dat.
\end{enumerate}

% \autoref přidá automaticky i název typu odkazu („oddíl“, „kapitola“).
% Pokud potřebujete český pád („v oddílu“), napište jej ručně s \ref.
Požadavky shrnuje \autoref{sec:pozadavky}. Na tuto analýzu
navazuje \autoref{chap:navrh}.

%% Copyright 2026 Petr Socha
%%
%% This work may be distributed and/or modified under the
%% conditions of the LaTeX Project Public License, either version 1.3c
%% of this license or (at your option) any later version.
%% The latest version of this license is in
%%   https://www.latex-project.org/lppl.txt
%% and version 1.3c is part of all distributions of LaTeX
%% version 2008 or later.
%%
%% This work has the LPPL maintenance status `maintained'.
%% The Current Maintainer of this work is Petr Socha.
%% Maintenance and bug reports: https://github.com/petrsocha/pslib-latex/issues
%%
%% The components of this work are listed in MANIFEST.txt.
%% The school logo is a separately licensed third-party asset.

% ============================================================================
% NÁVRH -- OBRÁZEK, TABULKA, POPISKY A ODKAZY
% ============================================================================
\chapter{Návrh}
\label{chap:navrh}

% Následující odstavce jsou pouze ilustrační výplňový text.
% Při psaní vlastní práce je nahraďte popisem navrženého řešení.
V~této kapitole je uvedena jednoduchá ukázka technického návrhu.
Text slouží především k~předvedení sazby obrázků, tabulek a~odkazů.
Ve skutečné práci by měl čtenář zjistit, jaké řešení bylo zvoleno,
proč je vhodné a~jak spolu jeho jednotlivé části souvisejí.

\section{Ukázkový obrázek}

Nejprve je vhodné popsat hlavní části navrhovaného systému a~jejich
vzájemné vazby. Přehledné schéma pomůže čtenáři pochopit princip řešení
dříve, než začne číst podrobnosti o~jednotlivých součástkách.
Přenos dat od senzoru znázorňuje \autoref{fig:tok-dat}:
mikrokontrolér předává naměřené hodnoty do počítače.

% Obrázky ukládejte do složky obrazky/ (PNG, JPG nebo PDF).
% Prostředí figure je plovoucí: [htbp] preferuje začátek/konec stránky,
% případně samostatnou stránku s obrázkem. Nenuťte obrázek na pevné místo.
% Popisek je POD obrázkem. Značku \label vložte ZA \caption.
% Odkaz v textu vložte pomocí \autoref{fig:tok-dat}.
\begin{figure}[htbp]
  \centering
  \includegraphics[width=0.78\textwidth]{obrazky/ukazkove-schema.pdf}
  \caption{Ukázkové schéma toku dat (vlastní zpracování)}
  \label{fig:tok-dat}
\end{figure}

Senzor získává měřenou veličinu, mikrokontrolér zpracuje přijaté údaje
a~počítač je zobrazí. Ve skutečné práci je vhodné připojit i~zdůvodnění
výběru jednotlivých součástek. Jejich parametry lze stručně porovnat
v~tabulce, která bývá přehlednější než souvislý výčet hodnot.

\section{Ukázková tabulka}

% Hodnoty jsou převzaty z oficiální dokumentace Arduino UNO R3.
% Tabulka má popisek NAD tabulkou a zarovnání pomocí prostředí tabular.
Dokumentace výrobce uvádí u~desky Arduino UNO R3 například procesor
ATmega328P a~takt 16\,MHz~\cite{arduinoUNO}.
Vybrané parametry shrnuje \autoref{tab:arduino-parametry}.

\begin{table}[htbp]
  \centering
  \caption{Vybrané parametry desky Arduino UNO R3 podle výrobce~\cite{arduinoUNO}}
  \label{tab:arduino-parametry}
  \begin{tabular}{ll}
    \hline
    Parametr & Hodnota \\
    \hline
    Mikrokontrolér & ATmega328P \\
    Taktovací frekvence & 16\,MHz \\
    Digitální vstupy/výstupy & 14 \\
    Analogové vstupy & 6 \\
    \hline
  \end{tabular}
\end{table}

Podrobnější parametry bývají uvedené v~datasheetu.
Jedním z~příkladů je technický list výrobce
STMicroelectronics~\cite{stm32datasheet}. Při porovnávání je nutné
ověřit jednotky a~konkrétní variantu součástky.

% U převzatých obrázků a údajů uveďte zdroj. U vlastních napište
% „vlastní zpracování“. Tabulky nepořizujte jako screenshoty.

%% Copyright 2026 Petr Socha
%%
%% This work may be distributed and/or modified under the
%% conditions of the LaTeX Project Public License, either version 1.3c
%% of this license or (at your option) any later version.
%% The latest version of this license is in
%%   https://www.latex-project.org/lppl.txt
%% and version 1.3c is part of all distributions of LaTeX
%% version 2008 or later.
%%
%% This work has the LPPL maintenance status `maintained'.
%% The Current Maintainer of this work is Petr Socha.
%% Maintenance and bug reports: https://github.com/petrsocha/pslib-latex/issues
%%
%% The components of this work are listed in MANIFEST.txt.
%% The school logo is a separately licensed third-party asset.

% ============================================================================
% REALIZACE -- ZDROJOVÝ KÓD A SPECIÁLNÍ ZNAKY
% ============================================================================
\chapter{Realizace}
\label{chap:realizace}

\section{Výpis zdrojového kódu}

% Krátký zdrojový kód vložte do prostředí lstlisting.
% Výpis dostane automatické číslo, popisek a volitelně čísla řádků.
% Diakritika v komentářích a řetězcích Pythonu je podporována.
\begin{lstlisting}[language=Python,caption={Výpočet průměru naměřených hodnot},label={lst:prumer}]
def prumer(hodnoty):
    if not hodnoty:
        return None
    return sum(hodnoty) / len(hodnoty)

# Český komentář i řetězec ve zdrojovém kódu
print("Průměrná hodnota:", prumer([20, 22, 24]))
\end{lstlisting}

Jednoduchou funkci v~jazyce Python ukazuje \autoref{lst:prumer}.
U~rozsáhlejších programů je vhodné uvést jen důležité úseky a~úplné
zdrojové soubory přiložit samostatně.

\section{Speciální znaky}

% Znaky %, _, &, # a $ maji v LaTeXu speciální význam.
% V běžném textu se zapisují jako \%, \_, \&, \# a \$.
Cesta k~souboru může vypadat jako \texttt{data\_2026.csv}.
Hodnota 50\,\% znamená polovinu celku; text \texttt{A\&B} obsahuje znak
ampersandu. Příkaz \verb|\texttt{...}| sází neproporcionálním písmem.

%% Copyright 2026 Petr Socha
%%
%% This work may be distributed and/or modified under the
%% conditions of the LaTeX Project Public License, either version 1.3c
%% of this license or (at your option) any later version.
%% The latest version of this license is in
%%   https://www.latex-project.org/lppl.txt
%% and version 1.3c is part of all distributions of LaTeX
%% version 2008 or later.
%%
%% This work has the LPPL maintenance status `maintained'.
%% The Current Maintainer of this work is Petr Socha.
%% Maintenance and bug reports: https://github.com/petrsocha/pslib-latex/issues
%%
%% The components of this work are listed in MANIFEST.txt.
%% The school logo is a separately licensed third-party asset.

% ============================================================================
% TESTOVÁNÍ -- VZORCE, MĚRNÉ JEDNOTKY A VÝSLEDKY
% ============================================================================
\chapter{Testování}
\label{chap:testovani}

\section{Matematický zápis}

V~odborné práci lze uvádět i~matematické vztahy. Například Ohmův zákon
vyjadřuje závislost proudu $I$ na napětí $U$ a~odporu $R$:

% Samostatný vzorec v prostředí equation se čísluje automaticky.
% Odkaz na vzorec funguje stejně jako odkaz na obrázek nebo tabulku.
\begin{equation}
  I = \frac{U}{R}
  \label{eq:ohm}
\end{equation}

Pro jednoduchý obvod podle \autoref{eq:ohm} s~$U=5\,\mathrm{V}$
a~$R=1\,\mathrm{k}\Omega$ vychází $I=5\,\mathrm{mA}$.

\section{Odborná přednáška jako zdroj}

% Skripta a přednáškové materiály z VŠ lze citovat stejně jako jiné zdroje.
% V literatura.bib je ukázka veřejně dostupné přednášky z MIT.
Vysokoškolské přednáškové poznámky mohou sloužit jako zdroj například
při vysvětlování filtrů. Materiál MIT \emph{Lecture 18: Filters}
popisuje mimo jiné dolní a~horní propust~\cite{mitfilters}.

U~jednoduché RC dolní propusti lze mezní frekvenci vypočítat ze vztahu
\begin{equation}
  f_{\mathrm c} = \frac{1}{2\pi RC}.
  \label{eq:rc}
\end{equation}

Při uvádění výsledků nestačí samotný výpočet; stručně vysvětlete,
co znamenají a~zda odpovídají stanoveným požadavkům.

%% Copyright 2026 Petr Socha
%%
%% This work may be distributed and/or modified under the
%% conditions of the LaTeX Project Public License, either version 1.3c
%% of this license or (at your option) any later version.
%% The latest version of this license is in
%%   https://www.latex-project.org/lppl.txt
%% and version 1.3c is part of all distributions of LaTeX
%% version 2008 or later.
%%
%% This work has the LPPL maintenance status `maintained'.
%% The Current Maintainer of this work is Petr Socha.
%% Maintenance and bug reports: https://github.com/petrsocha/pslib-latex/issues
%%
%% The components of this work are listed in MANIFEST.txt.
%% The school logo is a separately licensed third-party asset.

% ============================================================================
% ZÁVĚR -- ZHODNOCENÍ PRÁCE
% ============================================================================
% Zaver je nečíslovaný, stejne jako Uvod, ale zobrazi se v obsahu.
\necislo{Závěr}

V~závěru stručně uveďte, zda se podařilo splnit stanovené cíle,
co bylo výsledkem práce a~jaká omezení řešení má. Můžete rovněž naznačit
možnosti dalšího rozvoje.

% Do skutečné práce nepatří nové informace, které nebyly vysvětleny
% v předchozích kapitolách. Nepovažujte tento vzorový odstavec za zadání.

%TC:ignore

% ----------------------------------------------------------------------------
% 7. ZÁVĚREČNÉ ČÁSTI -- PODLE POTŘEBY ZAPNĚTE NEBO VYPNĚTE
% ----------------------------------------------------------------------------
% První závěrečná část začíná na pravé straně, další mohou navazovat
% bez prázdného rubu. Podle potřeby zakomentujte nepotřebné seznamy.
\zaverecnecasti

% Seznam zkratek: jeho položky zapisujte do souboru zkratky.tex.
\necislo{Seznam zkratek a odborných výrazů}
%% Copyright 2026 Petr Socha
%%
%% This work may be distributed and/or modified under the
%% conditions of the LaTeX Project Public License, either version 1.3c
%% of this license or (at your option) any later version.
%% The latest version of this license is in
%%   https://www.latex-project.org/lppl.txt
%% and version 1.3c is part of all distributions of LaTeX
%% version 2008 or later.
%%
%% This work has the LPPL maintenance status `maintained'.
%% The Current Maintainer of this work is Petr Socha.
%% Maintenance and bug reports: https://github.com/petrsocha/pslib-latex/issues
%%
%% The components of this work are listed in MANIFEST.txt.
%% The school logo is a separately licensed third-party asset.

% Zkratky zapisujte po jednom radku; zbytecne polozky muzete odstranit.
% Seznam by mel obsahovat jen zkratky skutecne pouzite v praci.
\begin{description}[style=nextline,leftmargin=2.8cm,labelwidth=2.5cm]
  \item[ADC] Analog-to-Digital Converter -- analogově-digitální převodník.
  \item[MCU] Microcontroller Unit -- mikrokontrolér.
  \item[RC] Obvod tvořený rezistorem (R) a kondenzátorem (C).
\end{description}


% Seznam obrázků se vytváří z příkazů \caption u obrázků.
\seznamobrazku

% Seznam tabulek je volitelný; zapněte jej odstraněním %.
%\seznamtabulek

% Použité zdroje se načítají z literatura.bib; v textu použijte \cite{...}.
\printbibliography[heading=bibintoc,title={Použité zdroje}]

% ----------------------------------------------------------------------------
% 8. PŘÍLOHY
% ----------------------------------------------------------------------------
% Další přílohy můžete vytvářet jako soubory ve složce prilohy/.
\zacatekpriloh
%% Copyright 2026 Petr Socha
%%
%% This work may be distributed and/or modified under the
%% conditions of the LaTeX Project Public License, either version 1.3c
%% of this license or (at your option) any later version.
%% The latest version of this license is in
%%   https://www.latex-project.org/lppl.txt
%% and version 1.3c is part of all distributions of LaTeX
%% version 2008 or later.
%%
%% This work has the LPPL maintenance status `maintained'.
%% The Current Maintainer of this work is Petr Socha.
%% Maintenance and bug reports: https://github.com/petrsocha/pslib-latex/issues
%%
%% The components of this work are listed in MANIFEST.txt.
%% The school logo is a separately licensed third-party asset.

\chapter{Seznam přiložených souborů}

% Výčet je pouze ukázkový: nahraďte názvy vlastními soubory.
% Odevzdává se výsledné PDF i CELÝ editovatelný LaTeXový projekt.
% Projekt obsahuje main.tex, všechny kapitoly, bibliografii, obrázky,
% potřebné soubory šablony a případné další součásti práce.

Na přiloženém datovém nosiči se nacházejí následující soubory a složky:

\begin{itemize}
  \item \textbf{\path{MP2026-Novak-Jan-P4B-Tepelne_cerpadlo.pdf}} -- tisknutelná verze dokumentace práce,
  \item \textbf{\path{MP2026-Novak-Jan-P4B-Tepelne_cerpadlo/}} -- editovatelná verze dokumentace práce; kompletní projekt v~LaTeXu včetně souborů \texttt{.tex}, bibliografie, obrázků a~dalších součástí potřebných k~překladu,
  \item \textbf{Výkresy} -- kompletní výkresová dokumentace,
  \item \textbf{Aplikace} -- zdrojové kódy.
\end{itemize}

% Název souboru PDF a složky projektu si přizpůsobte své práci.
% Nepřikládejte pouze main.tex: bez dalších souborů nemusí být projekt přeložitelný.


\end{document}
%TC:endignore
